% Originally: content/week-03.tex, \section{Compactness --- continued} (Corsin Week 3, Monday)

% Corsin taught this a week after the definition, and opened it as "Compactness --- continued"
% with a back-pointer. In a topic chapter the two sit together, so the heading names its content
% and the \continuedfrom pointer is gone.
\section{Heine--Borel}
\label{sec:heine_borel}

\begin{theorem}[Heine--Borel]
\label{thm:heine_borel}
A subset $K \subseteq \mathbb{R}^n$ equipped with the \textbf{standard metric} is
\textbf{compact} if and only if it is \newterm{closed} and \newterm{bounded}
\germanfor{beschr\"ankt}.
\end{theorem}

\begin{ainote}
Neither Corsin nor this document proves \cref{thm:heine_borel}; it is proved in the lecture, and
the full argument is longer than anything else in these two chapters. It is worth knowing what
goes into it, though, since every ingredient has already appeared:
\begin{itemize}
  \item \textbf{($\Rightarrow$)} is short and general. A compact metric space is bounded and
    complete (\cref{lem:sequentially_compact_totally_bounded},
    \cref{ex:compact_finite_complete}), and a complete subset of the complete space
    $\mathbb{R}^n$ is closed (\cref{rem:complete_vs_closed}).
  \item \textbf{($\Leftarrow$)} is where the real content sits, and it is exactly
    \textbf{Bolzano--Weierstrass}: a bounded sequence in $\mathbb{R}^n$ has a convergent
    subsequence, obtained by bisecting a box $n$ times over, or by extracting coordinatewise
    (\cref{prop:Rn_coordinatewise_complete}) $n$ times in succession. Closedness then puts the
    limit back inside $K$, giving sequential compactness, which is compactness by
    \cref{thm:sequential_topological_compactness_equivalence}.
\end{itemize}
The reason the theorem is special to $\mathbb{R}^n$ with the standard metric is that both steps
use the completeness of $\mathbb{R}$ and the finiteness of the dimension.
\Cref{ex:heine_borel_fails} below shows how thoroughly it fails otherwise.
\end{ainote}

% Source: Lukas Krause/Week 3 March 2-6 Notes Lukas Krause.pdf, p. 2
% Generator: Gemini 3.6 Flash (High)
\begin{proposition}[Topologically compact subsets of $\mathbb{R}^n$ are bounded]
\label{prop:topologically_compact_is_bounded}
If a subset $S \subseteq \mathbb{R}^n$ is topologically compact, then $S$ is bounded.
\end{proposition}

\begin{proof}
Consider the collection of open balls $\mathcal{U} := \{B_k(0) \mid k \in \mathbb{N}_{\ge 1}\}$ centered at the origin. Since every point $x \in \mathbb{R}^n$ has $|x| < k$ for some large integer $k$, $\mathcal{U}$ forms an open cover of $\mathbb{R}^n$, and hence an open cover of $S$:
\[
S \;\subseteq\; \mathbb{R}^n \;=\; \bigcup_{k=1}^\infty B_k(0).
\]
By topological compactness of $S$, there exists a finite subcover $\{B_{k_1}(0), B_{k_2}(0), \dots, B_{k_m}(0)\}$. Setting $K := \max\{k_1, k_2, \dots, k_m\}$, we obtain
\[
S \;\subseteq\; \bigcup_{j=1}^m B_{k_j}(0) \;=\; B_K(0).
\]
Therefore, $S$ is contained in the ball of radius $K$ centered at the origin, which means $S$ is bounded.
\end{proof}

% Supplement: Sascha Brack/Class Notes/Week_02_Notes_Friday_Updated.pdf, p. 7
\begin{example}[Compactness via Heine-Borel and continuity]
\label{ex:compactness_heine_borel_examples}
Consider the following subsets and determine whether they are compact under the standard Euclidean metric:
\begin{enumerate}[label=\textbf{(\alph*)}]
  \item \textbf{Is $\mathbb{N} \subset \mathbb{R}$ compact?} No. While it is closed in $\mathbb{R}$, it is not bounded. By Heine-Borel, it cannot be compact.
  \item \textbf{Is $[0,1] \times [0,1] \subset \mathbb{R}^2$ complete?} Yes. The square $[0,1]^2$ is closed and bounded in $\mathbb{R}^2$, hence compact by Heine-Borel. As every compact metric space is complete (\cref{ex:compact_finite_complete}), it is complete.
  \item Let $f: \mathbb{R} \to \mathbb{R}^3, x \mapsto (\sin x, \cos x, 1)$. \textbf{Is $f([0,\pi/2])$ compact in $\mathbb{R}^3$?} Yes. We do not need to show this manually using sequences or open covers. Since the domain $[0,\pi/2]$ is compact in $\mathbb{R}$ (closed and bounded) and the function $f$ is continuous, its image $f([0,\pi/2])$ is compact in $\mathbb{R}^3$ (by \cref{item:continuity_preserves_compact}).
\end{enumerate}
\end{example}

% Supplement: Sascha Brack/Class Notes/Week_03_Notes_Monday.pdf, p. 3
\begin{exercise}[Topological properties of various subsets]
\label{ex:topological_properties_table}
Decide whether the following subsets are open, closed, compact, or complete (under the standard Euclidean metric):
\begin{enumerate}[label=\textbf{(\alph*)}]
  \item $[0,1] \subseteq \mathbb{R}$
  \item $\mathbb{Q} \subseteq \mathbb{R}$
  \item $B_1(0) \subseteq \mathbb{R}^3$
  \item $B_1(0) \cap B_1(1) \subseteq \mathbb{R}^2$
  \item $\{0\} \cup \{1/n \mid n \in \mathbb{N}_{\ge 1}\} \subseteq \mathbb{R}$
  \item $\bigcup_{n=1}^\infty B_{1/n}(n) \subseteq \mathbb{R}$
  \item $(-\infty, 1] \subseteq \mathbb{R}$
  \item $f([0,1])$ where $f : \mathbb{R} \to \mathbb{R}^n$ is continuous
  \item $f^{-1}([0,1])$ where $f : \mathbb{R}^n \to \mathbb{R}$ is continuous
  \item $f^{-1}((0,1))$ where $f : \mathbb{R}^n \to \mathbb{R}$ is continuous
  \item $S := \operatorname{span}\{(1,0,0)\transp, (3,1,0)\transp\} \subseteq \mathbb{R}^3$
  \item $\{x \in \mathbb{R}^3 \mid \exists s \in S \text{ s.t.\ } d(x, s) < 1\}$ for a subset $S \subseteq \mathbb{R}^3$
\end{enumerate}
\end{exercise}

\begin{exercisesolution}[Proof of \cref{ex:topological_properties_table}]
\begin{enumerate}[label=\textbf{(\alph*)}]
  \item $[0,1] \subseteq \mathbb{R}$: \textbf{Closed}, \textbf{compact}, \textbf{complete}. (Not open).
  \item $\mathbb{Q} \subseteq \mathbb{R}$: \textbf{None}. It is neither open nor closed, not bounded (so not compact), and not complete.
  \item $B_1(0) \subseteq \mathbb{R}^3$: \textbf{Open}.
  \item $B_1(0) \cap B_1(1) \subseteq \mathbb{R}^2$: \textbf{Open} (finite intersection of open sets).
  \item $\{0\} \cup \{1/n \mid n \in \mathbb{N}_{\ge 1}\} \subseteq \mathbb{R}$: \textbf{Closed}, \textbf{compact}, \textbf{complete} (bounded and contains its only limit point $0$).
  \item $\bigcup_{n=1}^\infty B_{1/n}(n) \subseteq \mathbb{R}$: \textbf{Open} (arbitrary union of open sets).
  \item $(-\infty, 1] \subseteq \mathbb{R}$: \textbf{Closed}, \textbf{complete}.
  \item $f([0,1])$: \textbf{Closed}, \textbf{compact}, \textbf{complete} (continuous image of a compact set).
  \item $f^{-1}([0,1])$: \textbf{Closed}, \textbf{complete} (continuous preimage of a closed set).
  \item $f^{-1}((0,1))$: \textbf{Open} (continuous preimage of an open set).
  \item $S := \operatorname{span}\{(1,0,0)\transp, (3,1,0)\transp\} \subseteq \mathbb{R}^3$: \textbf{Closed}, \textbf{complete} (it is a $2$-dimensional subspace, hence closed in $\mathbb{R}^3$).
  \item $\{x \in \mathbb{R}^3 \mid \exists s \in S \text{ s.t.\ } d(x, s) < 1\}$: \textbf{Open}. This set is equivalently written as $\{x \in \mathbb{R}^3 \mid d(x, S) < 1\}$ or $\bigcup_{s \in S} B_1(s)$, which is an arbitrary union of open balls.
\end{enumerate}
\end{exercisesolution}

This captures our intuition of a \qt{compact} set. We will briefly show that Heine--Borel is
\textbf{false} for $\mathbb{R}^n$ with an \textbf{arbitrary metric}.

\begin{exercise}[Heine--Borel fails for arbitrary metrics]
\label{ex:heine_borel_fails}
Let $(X,d)$ be a metric space.
\begin{enumerate}[label=\textbf{(\alph*)}]
  \item \textit{(Optional)} Show that $\tilde d(x,y) := \dfrac{d(x,y)}{1 + d(x,y)}$ is a metric on
    $X$.
  \item Show that, if $U \subseteq X$ is open with respect to $d$, then $U$ is open with respect
    to $\tilde d$.
  \item Show that, for $X = \mathbb{R}$, $d(x,y) = \dfrac{|x-y|}{1+|x-y|}$,
    $\mathbb{N} \subseteq X$ is closed and bounded but not compact.
\end{enumerate}
\end{exercise}

% Correction: Claude Opus 5 (High) --- part (c) read $X = \mathbb{R}^n$ in the source. Flagged by
% Gemini 3.1 Pro (2026-08-11) as confusing to leave standing in the statement, and it is: the
% ainote below diagnosed the defect but the reader still met the wrong hypothesis first.
\begin{ainote}
Part \textbf{(c)} is stated in the source for $X = \mathbb{R}^n$, but it then works with
$\mathbb{N} \subseteq X$ and the absolute value $|x-y|$, which only parses for $n = 1$. It is
corrected to $X = \mathbb{R}$ above. Nothing is lost: the same argument runs verbatim in
$\mathbb{R}^n$ once $|x-y|$ is read as the Euclidean norm and $\mathbb{N}$ as
$\mathbb{N} \times \{0\}^{n-1}$.
\end{ainote}

\begin{aiexample}[Open Cover of $(0,1)$ with No Finite Subcover]
% Generator: Gemini 3.6 Flash (Medium)
Consider the open interval $X = (0,1)$ under the standard Euclidean metric. The family of open sets
\[
U_n := \left(\frac{1}{n}, 1\right) \quad \text{for } n \in \mathbb{N}_{\ge 2}
\]
forms an open cover of $(0,1)$ because for any $x \in (0,1)$, there exists $n \in \mathbb{N}$ with $1/n < x$. However, any finite subcollection $\{U_{n_1}, \dots, U_{n_k}\}$ has a maximum index $N := \max\{n_1, \dots, n_k\}$, so its union is $U_N = (1/N, 1)$, which fails to cover points in $(0, 1/N]$. Thus $(0,1)$ is bounded but not compact.
\end{aiexample}

% Supplement: Analysis_II_Script_v1.pdf, p. 22
\begin{example}[Heine--Borel needs $\mathbb{R}^n$ itself, not just the standard-looking metric]
\label{ex:rationals_not_compact}
$\mathbb{Q} \cap [0,2]$, with the metric inherited from the standard metric on $\mathbb{R}$, is
closed and bounded \emph{as a subset of $\mathbb{Q}$}, and yet it is not compact. The family
\[\mathbb{Q}\cap[0,2] \ =\ \big(\mathbb{Q}\cap[0,\sqrt2)\big) \ \cup\!\!
  \bigcup_{p\in\mathbb{Q},\,p>\sqrt2}\!\big(\mathbb{Q}\cap(p,2]\big)\]
is an open cover of $\mathbb{Q}\cap[0,2]$ (each set is the trace of an open subset of
$\mathbb{R}$, hence open in the subspace $\mathbb{Q}\cap[0,2]$): every rational in $[0,2]$ is
either $<\sqrt2$ or $>\sqrt2$, since $\sqrt2 \notin \mathbb{Q}$. Any finite subcover uses only
finitely many of the sets $\mathbb{Q}\cap(p,2]$, hence has a largest such $p$, call it
$p_{\max}$, and then misses every rational in $(\sqrt2,\,p_{\max}]$, of which there are infinitely
many. This is sharper than \cref{ex:compactness_heine_borel_examples}\textbf{(a)}'s $\mathbb{N}$.
There, compactness fails because the set is unbounded. Here it fails on a closed \emph{and}
bounded set, purely because the ambient space $\mathbb{Q}$ is not $\mathbb{R}^n$. The
irrationality of $\sqrt2$ is doing all the work, leaving a \qt{hole} exactly where the cover would
need to close up.
\end{example}

% Kept inline: solved directly in Corsin Nick/Class Notes/Week 3.pdf, pp. 1--2.
\begin{exercisesolution}[Proof of \cref{ex:heine_borel_fails}]
\begin{enumerate}[label=\textbf{(\alph*)}]
  \item Non-negativity, definiteness and symmetry are inherited directly from the
    corresponding properties of $d$, since $t \mapsto \tfrac{t}{1+t}$ is non-negative on
    $[0,\infty)$ and vanishes only at $t=0$.

    \textit{Triangle inequality.} We want to show
    \[
    \frac{d(x,y)}{1+d(x,y)} \leq \frac{d(x,z)}{1+d(x,z)} + \frac{d(z,y)}{1+d(z,y)}.
    \]
    Denote $d(x,y) = d_1$, $d(x,z) = d_2$, $d(z,y) = d_3$ and multiply out:
    \[
    \begin{aligned}
    &d_1(1+d_2)(1+d_3) \leq d_2(1+d_1)(1+d_3) + d_3(1+d_1)(1+d_2) \\
    \iff\ & (d_1 + d_1d_2)(1+d_3) \leq (d_2 + d_2d_1)(1+d_3) + (d_3 + d_1d_3)(1+d_2) \\
    \iff\ & d_1 + d_1d_3 + d_1d_2d_3 + d_1d_2 \leq d_2 + d_3d_2 + d_2d_1d_3 + d_2d_1 \\
    & \qquad\qquad\qquad\qquad\qquad\quad + d_3 + d_3d_2 + d_1d_3 + d_1d_3d_2 \\
    \iff\ & d_1 \leq d_2 + 2d_2d_3 + d_3 + d_1d_2d_3 \\
    \Longleftarrow\ & 0 \leq 2d_2d_3 + d_1d_2d_3 \quad \text{together with } d_1 \leq d_2+d_3
      \text{ (the triangle inequality for } d\text{)}
    \end{aligned}
    \]
    which is true by non-negativity.
  \item Let $U \subseteq X$ be open. Then, given $x \in U$, we find $r > 0$ such that
    $B_r(x) := \{y \in X : d(x,y) < r\} \subseteq U$. Let $\varepsilon := \dfrac{r}{1+r}$. Then
    $\tilde B_\varepsilon(x) \subseteq B_r(x) \subseteq U$, where
    $\tilde B_\varepsilon(x) := \{y \in X : \tilde d(x,y) < \varepsilon\}$.
  \item Note that for all $x, y \in \mathbb{R}$, $\dfrac{|x-y|}{1+|x-y|} < 1$. So
    $\mathbb{N} \subseteq B_1(0)$, thus it is \textbf{bounded} with respect to this metric. Also,
    $\mathbb{N}$ is \textbf{closed}, as follows from \textbf{(b)}. But $x_n = n$ defines a
    sequence in $\mathbb{N}$ with no convergent subsequence, thus $\mathbb{N}$ is \textbf{not
    compact}.
\end{enumerate}
\end{exercisesolution}

\begin{ainote}
The cancellation in the derivation above originally contained a stray $d_2d_1d_2$ where the
expansion gives $d_2d_1d_3$; the cancelled terms and the conclusion $0 \leq 2d_2d_3 + d_1d_2d_3$
are unaffected. Re-derived cleanly above.
\end{ainote}

% Generator: Claude Opus 5 (Medium)
\begin{remark}[The moral: boundedness is not a topological property]
\label{rem:boundedness_not_topological}
It is worth saying out loud what \cref{ex:heine_borel_fails} has just demonstrated, because the
three parts are more striking together than separately.

Part \textbf{(b)} shows that $d$ and $\tilde d$ have the \emph{same open sets}. (Only one
direction is asked for, but the converse follows the same way: given
$\tilde B_\varepsilon(x)$, the choice $r := \tfrac{\varepsilon}{1-\varepsilon}$ gives
$B_r(x) \subseteq \tilde B_\varepsilon(x)$.) Two metrics with the same open sets have the same
convergent sequences, the same continuous functions, and, crucially, the same \textbf{compact}
sets, since compactness is defined purely by open covers (\cref{def:compact_metric_space}).

Part \textbf{(c)} shows that they do \emph{not} have the same bounded sets. Indeed
$\tilde d < 1$ always, so under $\tilde d$ \emph{every} subset of \emph{every} metric space is
bounded. The notion collapses entirely.

Putting those together: \textbf{compactness is topological, boundedness is not.} Boundedness is
a property of the chosen metric, not of the open sets it generates, so no theorem of the form
\qt{compact if and only if closed and bounded} can hold for a general metric. That is precisely why
\Cref{thm:heine_borel} carries the hypothesis \textbf{standard metric}. That hypothesis is not a
technicality to be skimmed; it is the entire content of the theorem.

% Reclassified from ainote to remark: it introduces terminology with \newterm and separates two
% notions by an explicit example, so it is mathematics.
% Feedback: Gemini 3.1 Pro (2026-08-11) --- singled out for placing total boundedness against
% boundedness at exactly the point where the reader has just watched boundedness collapse. Keep
% it nested here; moving it to the definition of total boundedness would lose the contrast.
\begin{remark}[What boundedness should have been]
The right general replacement is already available: compact if and only if complete \emph{and}
\newterm{totally bounded} (\cref{def:totally_bounded},
\cref{rem:compact_iff_complete_totally_bounded}). Total boundedness is the notion that boundedness
was trying and failing to be, and $\mathbb{N}$ under $\tilde d$ separates the two cleanly. For
$\varepsilon < \tfrac12$ the ball $\tilde B_\varepsilon(n)$ contains no integer other than $n$
itself, so no finite family of $\varepsilon$-balls covers $\mathbb{N}$. The set is therefore
bounded but not totally bounded, and that is exactly the gap through which compactness escapes.
\end{remark}
\end{remark}

% Source: old_exams/august2025.pdf (Prof. Laura Kobel-Keller), Analysis II examination,
% August 2025, Frage 10, printed p. 9
% Extractor: Claude Opus 5 (High)
\begin{exercise}[True or False \textnormal{(what actually characterises compactness)}]
\label{ex:kobel_compactness_characterisation}
Let $E \subseteq X$ be a subset of a metric space. For each of the following completions of the
statement \qt{$E$ is compact if \dots}, decide whether it is true or false. More than one may be
correct, and so may none.
\begin{enumerate}[label=\textbf{(\alph*)}]
  \item $E$ is closed and bounded.
  \item $E$ is simultaneously open and closed.
  \item $E$ is complete and totally bounded.
  \item $E$ is open and totally bounded.
\end{enumerate}
\exinfo{This exercise is taken from the Analysis II examination of August 2025
(Prof.\ Laura Kobel-Keller), Question 10, the opening question of the multiple-choice section.}
\end{exercise}
